\begin{center}
    
{\fontsize{16pt}{19pt}\selectfont \textbf{\MakeUppercase{Integrasi BERTopic dan Pre-trained Language Model (IndoBERT) untuk Pemetaan Aspect-Based Sentiment Analysis pada Destinasi Wisata alam dan religi Bangkalan}}}

Oleh \\[3pt]
    Seinal Arifin \\[8pt] % Jarak antara nama dan NIM
    Nim: 23041100149
\end{center}
\vspace{2cm}
\begin{center}
  \textbf{\large Abstrak}
  \addcontentsline{toc}{section}{Abstrak}
\end{center}
\noindent
Pariwisata di Kabupaten Bangkalan memiliki potensi perkembangan yang sangat besar, namun pihak pengelola destinasi seringkali mengalami kesulitan melakukan evaluasi tepat sasaran karena kurangnya analisis mendalam terhadap ulasan dari pengunjung. Mengetahui kecenderungan opini masyarakat secara umum tidaklah cukup; diperlukan pemetaan sentimen positif dan negatif pada setiap aspek atau topik wisata yang lebih spesifik. Penelitian ini mengusulkan solusi berupa pembangunan sistem analisis sentimen berbasis aspek menggunakan algoritma BERTopic untuk proses ekstraksi topik dan model IndoBERT untuk klasifikasi opini. Pengujian keandalan sistem dilakukan melalui tahapan pembersihan data teks dengan penyesuaian stopword khusus, prediksi label sentimen, serta klasterisasi topik secara dinamis menggunakan HDBSCAN terhadap ribuan ulasan pariwisata. Hasil penelitian menunjukkan bahwa pemodelan menggunakan IndoBERT berhasil mencapai tingkat akurasi sebesar 83,7\% dalam mengklasifikasikan sentimen. Selanjutnya, algoritma ekstraksi berhasil mengelompokkan ulasan secara otomatis ke dalam 15 topik aspek pariwisata, sekaligus memetakan aspek baik maupun aspek negatif di dalamnya. Seluruh hasil analisis divisualisasikan ke dalam sebuah dasbor interaktif guna mendukung pihak terkait melakukan evaluasi destinasi secara terukur dan berkala. \\
\noindent
\textbf{Kata kunci} : Pariwisata, Analisis Sentimen, Aspek, BERTopic, IndoBERT.